\documentclass[10pt,a4paper]{article}
\usepackage{amsmath,amssymb,amsthm,graphicx}
\pagestyle{plain}
\begin{document}
\section{Symmetries of the square}
The dihedral group of the square has eight elements. In the picture, the shaded cell is the reference position.
\begin{figure}[t]
\centering
\includegraphics[height=2cm]{images/square.png}\quad
\includegraphics[scale=2]{images/square.png}
\caption{The reference square at two sizes.}
\label{fig:two}
\end{figure}

Figure~\ref{fig:two} appears above; the group acts on it by rotations by multiples of $90^\circ$ and four reflections.

\begin{table}[b]
\centering
\begin{tabular}{@{}lcc@{}}
\hline
Type & Count & Fixed points\\ \hline
Identity & 1 & 4 corners\\
Rotation by $90^\circ$ & 2 & none\\
Rotation by $180^\circ$ & 1 & none\\
Reflection & 4 & 2 corners or 0\\ \hline
\end{tabular}
\caption{Conjugacy-type summary.}
\label{tab:sym}
\end{table}

By Burnside's lemma, the number of orbits of colourings of the four corners with two colours is
\[
\frac18\bigl(2^4+2\cdot2+2^2+2\cdot2^3+2\cdot2^2\bigr)=\frac{16+4+4+16+8}{8}=6.
\]
Table~\ref{tab:sym} supports the count.
\end{document}
